\subsection{From analytic irreducibility to a single finite branch}
The following classical equivalence identifies the local domains to which
the conductor calculation applies.  We give the passage explicitly because
finite normalization by itself does not exclude several branches.
The forward locality implication is also the proposition of
\cite[\S2, p.~49, proof pp.~49--50]{CahenGrazziniHaouat1982}.

\begin{proposition}[Analytic irreducibility and finite local normalization]
\label{prop:analytic-one-branch}
Let $(R,\mathfrak m)$ be a one-dimensional Noetherian local domain,
let $K$ be its fraction field, and let $V$ be its integral closure in $K$.
Then $\widehat R$ is a domain if and only if $V$ is a finite $R$-module
and is local.  In that case $V$ is a discrete valuation ring and
$\pi^mV\subset R$ for some uniformizer $\pi$ and integer $m\geq1$.
\end{proposition}
\begin{proof}
Suppose first that $\widehat R$ is a domain.  It is reduced, so $R$
is analytically unramified.  The normalization-finiteness theorem
\cite[Lemma 10.162.10(4), tag 032Y]{StacksCompletionAnalytic} makes $V$
a finite $R$-module.  Completion is flat, and hence the inclusion
$V\subset K$ gives an injection
\[
 V\otimes_R\widehat R\ \hookrightarrow\
 K\otimes_R\widehat R
 =(R\setminus\{0\})^{-1}\widehat R.
\]
The localization on the right is a domain: faithful flatness keeps
every nonzero element of $R$ nonzero in $\widehat R$.  Thus
$B=V\otimes_R\widehat R$ is a domain.  It is finite over the complete
Noetherian local ring $\widehat R$, so it is a finite product of
complete local rings \cite[Lemma 10.160.2, tag 0325]{StacksCompletionAnalytic}.
A domain has no nontrivial idempotents, and therefore this product
has exactly one factor.  It follows that $B$ is local.

Reduction modulo $\mathfrak m\widehat R$ identifies
$B/\mathfrak mB$ with $V/\mathfrak mV$.  Every maximal ideal of the
finite integral $R$-algebra $V$ lies above $\mathfrak m$, so its
maximal ideals correspond to those of $V/\mathfrak mV$.  This quotient
is local because $B$ is local.  Hence $V$ is local.  It is Noetherian
by finiteness, normal by definition, and one-dimensional by integrality.
The discrete-valuation criterion
\cite[Lemma 10.119.7, tag 00PD]{StacksNormalDomains} now shows that $V$
is a discrete valuation ring.

Conversely, suppose that $V$ is finite and local.  The same criterion
makes it a discrete valuation ring.  Its maximal ideal $\mathfrak n$
is the radical of $\mathfrak mV$, so the $\mathfrak mV$-adic and
$\mathfrak n$-adic topologies are equivalent.  Exactness of completion
on finite modules sends $R\hookrightarrow V$ to
$\widehat R\hookrightarrow\widehat V$.  Here $\widehat V$ is the
completion of a discrete valuation ring in its maximal ideal, and is
again a discrete valuation ring.  In particular it is a domain, and
its subring $\widehat R$ is a domain.

Finally, choose finitely many $R$-module generators for $V$ in $K$.
Clearing their denominators gives $0\ne d\in R$ with $dV\subset R$.
In the discrete valuation ring, $dV=\pi^aV$ for some $a\geq0$.
Taking $m\geq\max(1,a)$ gives the asserted conductor inclusion.
\end{proof}

\begin{proof}[Proof of Theorem~\ref{thm:int-analytic-free}]
If $R/\mathfrak m$ is infinite, let $f\in\operatorname{Int}(R)$ have
degree at most $n$.  Choose $n+1$ elements of $R$ with distinct
residues.  Their Vandermonde determinant is a unit.  Solving for the
coefficients of $f$ from its values at these elements puts every
coefficient in $R$.  Therefore $\operatorname{Int}(R)=R[X]$ in this case.

Suppose that $R/\mathfrak m$ is finite.  By
Proposition~\ref{prop:analytic-one-branch}, its normalization $V$ is
a finite local discrete valuation ring and contains a power of its
maximal ideal inside $R$.  The residue field of $V$ is finite:
it is a quotient of the finite-dimensional $R/\mathfrak m$-algebra
$V/\mathfrak mV$.  Theorem~\ref{result-29} therefore proves freeness
of $\operatorname{Int}(R)$, by projective patching followed by
Kaplansky's theorem.

In both cases the rank is countably infinite.  In the finite-residue
case the patching construction makes the module a direct summand of
a countably generated free module, and hence countably generated.
Moreover
\[
 K\otimes_R\operatorname{Int}(R)=K[X],
\]
because $R[X]\subset\operatorname{Int}(R)\subset K[X]$.
The dimension of this vector space over $K$ is countably infinite.
The tensor and compatible plethory conclusions follow from
Theorem~\ref{result-24}, applied to the fraction-field extension.
\end{proof}
